\documentclass[10pt,a4paper]{amsart}
\usepackage[margin=19mm]{geometry}
\usepackage{amsmath,amssymb,mathtools}
\usepackage[scr=rsfs,cal=euler]{mathalfa}
\usepackage{xfrac}
\usepackage[unicode,hidelinks,bookmarks=false]{hyperref}
\newcommand{\ie}{i.e.\ }
\newcommand{\cf}{cf.\ }
\newcommand{\resp}{respectively\ }
\newcommand{\Subsection}{\subsection}
\providecommand{\nobreakdash}{\nobreak\hbox{-}}
\setlength{\emergencystretch}{2em}
\multlinegap0pt
\usepackage{mathtools}
\newtheorem{lemma}{Lemma}
\newtheorem{theorem}{Theorem}
\newcommand{\HJ}{\mathrm{HJ}}
\renewcommand{\ge}{\geqslant}
\renewcommand{\le}{\leqslant}
\DeclareMathOperator{\type}{type}
\title{Lower Bounds for the Hales--Jewett Numbers\\ via Symmetric and One-Weight Colorings}
\author{Younes Mouhib}
\date{Source: arXiv:2606.22155v2 (CC BY 4.0)}
\begin{document}
\maketitle

\noindent\emph{Source and licence.} Lower Bounds for the Hales--Jewett Numbers\\ via Symmetric and One-Weight Colorings. Version arXiv:2606.22155v2 (CC BY 4.0), distributed by arXiv under the Creative Commons Attribution 4.0 International licence, http://creativecommons.org/licenses/by/4.0/, as recorded in the arXiv licence metadata for this exact version and shown on the abstract page https://arxiv.org/abs/2606.22155v2. Full licence notice: \texttt{licenses/arxiv-2606.22155-CC-BY-4.0.txt}.\par\smallskip

\noindent\textit{Editorial scope.} Counted unit: the article's theorem thm:radix (source0.tex lines 617--642). The article's complete original proof of it is delivered with it (source0.tex lines 644--664, one \texttt{proof} environment of 21 lines). No mathematics is written, repaired or re-derived: the statement and the proof are the article's own lines, in the article's own notation, and no retained line is substituted or re-flowed.  Retained uncounted as the source-local context the proof needs: lines 405--416.  Nothing was omitted from any retained span, so the package carries no omission record.  No retained line contains a citation, so the wrapper carries no bibliography and no bibliography entry.  No retained line contains a figure environment, an image-inclusion command or drawing code, so no image is delivered and the package declares no asset dependency.  Editorial formatting only, in addition: an AMS \texttt{amsart} preamble that restates the article's own declarations and macros used by the retained lines, namely the environments \texttt{lemma}, \texttt{theorem} on separate counters, together with the standard packages those lines need.  The retained environments print in this document as Lemma 1, Theorem 1 in order of appearance, whereas the article prints the same items with the numbering of its own full document.  The complete native source of the article as distributed in the arXiv e-print for this exact version is bundled unchanged as \texttt{sources/203328/source0.tex}.
\begin{lemma}[Symmetric reduction]\label{lem:sym-reduction}
Let $\bar c\colon T_{n}^{(t)}\to[r]$ and let $c=\bar c\circ\type$. For
any root $\tau$ with $k$ active coordinates and an inactive type $v$,
and for every $a\in[t]$,
\begin{equation}\label{eq:type-identity}
\type\bigl(\tau(a)\bigr)=v+k\,e_{a}.
\end{equation}
Consequently, $c$ has no monochromatic line with at most $K$ active
coordinates if and only if $\bar c$ has no monochromatic corner tuple
$C_{k,v}$ with $1\le k\le K$. At $K=n$: $c$ is line-free if and only if
no corner tuple is monochromatic.
\end{lemma}

\begin{theorem}[{One-weight $=$ symmetric; the rank hierarchy is flat}]\label{thm:radix}
Fix $t\ge 2$, $r\ge 2$, $n\ge 1$. The colorings of $[t]^{n}$ of the
form $c_{\omega,\chi}$ are exactly the symmetric colorings.
Specifically, for the \emph{radix weight}
\[
\omega^{\ast}=(0,1,p,p^{2},\dots,p^{t-2})
\qquad\text{with an arbitrary integer } p>n
\]
\textup{(}primality plays no role\textup{)}, the functional
$\langle\omega^{\ast},\cdot\rangle$ is injective on $T_{n}^{(t)}$ and
carries each corner tuple $C_{k,v}$ bijectively onto a carry-free
lattice corner; consequently every symmetric coloring $\bar c$ is
realized as $c_{\omega^{\ast},\chi}$ with
$\chi(\langle\omega^{\ast},v\rangle)=\bar c(v)$, and
$c_{\omega^{\ast},\chi}$ is line-free if and only if no corner tuple is
monochromatic under $\bar c$ (Lemma~\ref{lem:sym-reduction}). Moreover
any coloring built from finitely many weights
\[
w\ \longmapsto\ \Psi\bigl(\langle\omega^{(1)},\type(w)\rangle,\dots,
\langle\omega^{(s)},\type(w)\rangle\bigr)
\]
is symmetric, hence one-weight: there is no multi-weight tier, and
\[
\HJ_{\mathrm{ow}}(t,r)=\HJ_{\mathrm{sym}}(t,r).
\]
\end{theorem}

\begin{proof}
On $T_{n}^{(t)}$ each entry satisfies $0\le v_{a}\le n<p$ for $a\ge 2$,
so $\langle\omega^{\ast},v\rangle=\sum_{a\ge 2}v_{a}p^{a-2}$ is the
base-$p$ integer with digit string $(v_{2},\dots,v_{t})$; it is
injective because $v_{1}$ is determined by
$v_{1}=n-\sum_{a\ge 2}v_{a}$. A corner point $v+ke_{a}$ with $a\ge 2$
has $a$-digit $v_{a}+k\le n<p$ (as $v_{a}\le n-k$), so no carry occurs
and $\langle\omega^{\ast},\cdot\rangle$ maps $C_{k,v}$ bijectively onto
the digit corner $\{(v_{2},\dots,v_{t})+k\hat e_{a}\}_{a\in[t]}$, where
$\hat e_{a}$ is the $a$-th standard basis vector of the digit space for
$a\ge 2$ and $\hat e_{1}:=0$ (the point $v+ke_{1}$ keeps the digit
string of $v$, since $\omega^{\ast}_{1}=0$). Hence
$\chi(\langle\omega^{\ast},v\rangle):=\bar c(v)$ is well defined
(extended arbitrarily off the image),
$c_{\omega^{\ast},\chi}=\bar c\circ\type$, and line-freeness transfers
by Lemma~\ref{lem:sym-reduction}. Every $c_{\omega,\chi}$ is symmetric;
a multi-weight coloring factors through $\type$, hence is symmetric,
hence one-weight by the above. Since the two classes coincide in every
dimension, the defining conditions of $\HJ_{\mathrm{ow}}$ and
$\HJ_{\mathrm{sym}}$ agree for each $n$, so the thresholds are equal.
\end{proof}

\end{document}
